% Generated draft: requires mathematical and source review.
\documentclass[11pt]{article}
\usepackage[margin=25mm]{geometry}
\usepackage{fontspec}
\setmainfont{Latin Modern Roman}
\usepackage{amsmath,amssymb,mathtools,unicode-math}
\setmathfont{Latin Modern Math}
\usepackage{enumitem,needspace,xurl,hyperref}
\setlength{\parindent}{0pt}
\setlength{\parskip}{5pt}
\setlength{\emergencystretch}{3em}
\newcommand{\N}{\mathbb{N}}
\newcommand{\Z}{\mathbb{Z}}
\newcommand{\Q}{\mathbb{Q}}
\newcommand{\R}{\mathbb{R}}
\newcommand{\C}{\mathbb{C}}
\newcommand{\F}{\mathbb{F}}
\newcommand{\A}{\mathbb{A}}
\newcommand{\PP}{\mathbb{P}}
\newcommand{\E}{\mathbb{E}}
\newcommand{\eps}{\varepsilon}
\newcommand{\dd}{\,\mathrm{d}}
\begin{document}
\section*{UnsolvedMath ID: 30000940}
\textbf{Source problem code:} catalogue entry
\subsection{Definitions and mathematical statement}
Let \(\zeta(s)\) be the Riemann zeta function, and assume the Riemann Hypothesis: every nontrivial zero of \(\zeta\) has real part \(1/2\). Let \(\rho'=\beta'+i\gamma'\) run over the nontrivial zeros of \(\zeta'(s)\) with positive ordinates tending to infinity. Let \(\gamma\) run over the positive ordinates of the nontrivial zeros of \(\zeta(s)\), and let \(\gamma^+\) denote the next such ordinate in the source's ordering.

Soundararajan's conjecture, as supplied, asks for the equivalence
\[
\liminf_{\substack{\zeta'(\beta'+i\gamma')=0\\\gamma'\to\infty}}
(\beta'-\tfrac12)\log\gamma'=0
\quad\Longleftrightarrow\quad
\liminf_{\gamma\to\infty}(\gamma^+-\gamma)\log\gamma=0.
\]
The left limit is along the designated derivative zeros, not over all real heights. In a customary ordering \(\gamma_1\leq\gamma_2\leq\cdots\) with multiplicity, the right side is
\[
\liminf_{n\to\infty}(\gamma_{n+1}-\gamma_n)\log\gamma_n=0.
\]
The extraction does not expressly specify multiplicity versus strictly distinct ordinates; this must be checked before fixing the successor convention.

The source additionally defines \(\gamma_c\) as the ordinate of a zeta zero nearest to \(\gamma'\), meaning one minimizing \(|\gamma_c-\gamma'|\). It supplies no tie-breaking convention, and \(\gamma_c\) does not occur in the displayed equivalence. No additional conclusion involving it is inferred. The intended nontrivial-zero convention for \(\zeta'\) should likewise be confirmed; derivative zeros are not silently restricted to \(0<\beta'<1\).
\subsection{Short English statement}
Assuming RH, do derivative zeros approach the critical line on the logarithmic scale exactly when consecutive zeta-zero ordinates have arbitrarily small logarithmically normalized gaps? This is the supplied Soundararajan equivalence.
\subsection{Sources}
\begin{itemize}
\item Local UnsolvedMath registry, permanent ID \(30000940\), source code catalogue entry, ``Soundararajan Conjecture on Zeta Zero Spacing.'' The original statement attributes the conjecture to Soundararajan; no individual proposer field is populated.
\item \emph{Analytic Number Theory} (2008), Oberwolfach Reports, DOI 10.4171/owr/2008/14, pp. 669--746, also listed as ``Oberwolfach Report 14/2008'': \url{https://doi.org/10.4171/owr/2008/14}.
\item Supplied reference, ``Garaev and Yildirim, zeros of zeta and zeta prime'': \url{https://doi.org/10.1093/imrn/rnm091}.
\item Supplied literature triage, checked 2026-08-21: status ``partially solved.'' It reports quantitative implications by Garaev and Yildirim, calls the stronger Soundararajan equivalence and displayed limiting assertions not fully settled, and describes those implications as involving ``prime gaps.'' The displayed question instead concerns gaps between zeta-zero ordinates; that terminology discrepancy is retained for review. This dated assessment derives from supplied OWR curation; time-sensitive status and cited claims require independent verification.
\end{itemize}
\end{document}
